\subsubsection{Future directions}
Recently, impressive advances in different tasks in Artificial Intelligence have been achieved using deep learning techniques.
Embedding-based language models, such as Word2Vec~\cite{mikolov1,mikolov2} and GloVe~\cite{glove}, have helped to improve performance in many NLP tasks. 
%Deep learning techniques are being explored with promising results for different problems. 
One of the most interesting and the most promising future directions for semantic parsing and question answering is further exploration of deep learning techniques in their context.

%The essential idea behind deep learning is to let the learning algorithm learn to extract its own features.
% Where traditional learning techniques relied on carefully engineered feature extractors, deep learning leaves the feature extraction up to the machine, with the option of pretraining feature vectors on large amounts of unsupervised data.
% The two main bottlenecks for such approaches, large amounts of data and computational power, are becoming less and less of an issue.

Using deep learning to better understand questions can be done by using (possibly custom-trained) word (, word sense and entity) embeddings, which capture their syntactic and semantic properties, as features to improve existing workflows.
However, a ``deeper'' approach would be to also devise new models that provide the machine with more freedom to figure out how to accomplish the task.
An excellent and very recent example in NLP is the Dynamic Memory Network (DMN~\cite{dmn}), that does not use any manually engineered features or problem-tailored models, and yet achieves state-of-the-art performance on all tested tasks, which are disjoint enough to leave one impressed (POS tagging, co-reference resolution, sentiment analysis and Question Answering on the bAbI dataset).
The DMN is one of the works focusing on attention and memory in deep learning that enables the neural network to reason more freely.
We share the belief that the investigation and application of more advanced deep learning models (such as DMN and NTM~\cite{ntm}) could yield impressive results for different tasks in AI, including question answering.

Recursive, convolutional (CNN) and recurrent (RNN) neural networks are widely used in recent neural network-based approaches.
Convolutional Neural Networks (CNN), a special case of recursive NNs are well-explored for computer vision.
Recursive NNs have also been applied for parsing and sentiment analysis. 
RNNs produce state-of-the-art results in speech processing as well as in NLP because of their natural vigor for processing variable-length sequences.
They have been applied for machine translation (SMT)~\cite{smtrnn}, language generation (NLG)~\cite{nlgrnn}, language modeling and more and are also fundamental for the success of the DMN and the NTM.

Even though the DMN has not yet been applied to our task of structured QA, some recent works, such as the relatively simple embedding-based work of Bordes et al. ~\cite{bordes2014question} (which outperformed ParaSempre~\cite{parasempre} on \textsc{WebQuestions}) and the SMT-like SP approach of~\cite{dong2016language} seem to acknowledge the promise of neural approaches with embeddings.

An additional interesting direction is the investigation of joint models for the sub-problems involved in question interpretation (EL, co-reference resolution, parsing, \dots).
Many tasks in NLP depend on each other to some degree, motivating the investigation of efficient approaches to make the decisions for those tasks jointly.
For example, co-reference resolution and EL can benefit from each other as entity information from a KB can serve as quite powerful features for co-reference resolution and co-reference resolution in turn can improve EL as it transfers KB features to phrases where anaphora refer to entities.
Factor Graphs (and Markov Networks) are by nature very well-suited for explicit joint models (e.g.~\cite{jointriedel}).
However, a more internal kind of joint inference could also be achieved within a neural architecture (e.g. the DMN).%, as exemplified by the ability of the DMN to resolve  %For example, the DMN mentioned above needs to resolve anaphora in order to correctly answer certain questions.
%Based on the results of the DMN on that dataset, it can be assumed that it actually learned to perform coreference resolution where needed without any explicit model or a supervision signal for coreference resolution.

However, it is worth noting that training advanced neural models and explicit joint models can be a difficult task because of the large number of training parameters and co-dependence of these parameters. Deep learning typically relies on the availability of large datasets. However, the whole task to be solved can be divided in two parts, one focusing on representation learning, which can accomplished in an unsupervised setting (with large amounts of data) and the second part relying on and possibly fine-tuning the representations obtained in the first part in a supervised training setting (requiring annotated task-specific data). For explicit joint models, data capturing the dependence between different task-specific parts of the models (e.g. annotated for both EL and co-reference) are required and the efficient training of such models is a very relevant current topic of investigation.

The concluding thought is that the further investigation of language~\cite{mikolov1,mikolov2,glove,sensembed} and knowledge modeling~\cite{tatec,transe,riedel2013,ntn,trescal} and powerful deep neural architectures with self-regulating abilities (attention, memory) as well as implicit or explicit joint models will continue to push the state of the art in QA.
Well-designed deep neural architectures, given proper supervision and powerful input models, have the potential to learn to solve many different NLU problems robustly with minimal customizations, eliminating the need for carefully engineered features, strict formalisms to extract complex structures or pipelines arranging problem-tailored algorithms.
We believe that these lines of research in QA could be the next yellow brick~\cite{yellowbrick} in the road to true AI, which has fascinated humanity since the ancient tales of Talos and Yan Shi's mechanical men.

% Local Variables:
% mode: LaTeX
% TeX-master: "paper"
% End:
%  LocalWords:  Vec GloVe NLP DMN POS bAbI NTM convolutional RNN NNs
%  LocalWords:  RNNs SMT NLG QA Bordes ParaSempre WebQuestions EL NLU
%  LocalWords:  anaphora Talos Yan Shi's
